% 12_formato_json.tex - Capítulo 12: Formato JSON y Validación
\chapter{Formato JSON de Escenarios y Validación Atómica}
\label{ch:json}

\section{Contrato de Interoperabilidad}
\label{sec:contrato-json}

\screenshotfigure{json-escenario}{0.90}{Archivo JSON de escenario utilizado por la aplicación.}

\screenshotfigure{json-validacion}{0.85}{Validación de un escenario JSON.}

El formato JSON es el \textbf{contrato de interoperabilidad} (Sección 7 del enunciado):
cualquier archivo válido que respete el formato debe cargarse \textbf{sin modificar el código
fuente}. El profesor entregará al menos un escenario no visto en la sustentación.

\section{Campos Obligatorios (RF-35)}
\label{sec:campos-obligatorios}

\begin{longtable}{p{3cm}p{3cm}p{8cm}}
\toprule
\textbf{Campo} & \textbf{Tipo} & \textbf{Descripción} \\
\midrule
\endfirsthead
\toprule
\textbf{Campo} & \textbf{Tipo} & \textbf{Descripción} \\
\midrule
\endhead
\midrule
\multicolumn{3}{r}{\textit{Continúa}} \\
\endfoot
\bottomrule
\endlastfoot
\code{version} & string & Versión del formato (ej. \code{"1.0"}) \\
\code{mapa.filas} & int $>0$ & Filas del grid \\
\code{mapa.columnas} & int $>0$ & Columnas del grid \\
\code{tipos\_terreno} & object & Catálogo: \code{nombre $\to$ \{costo, transitable\}} \\
\code{terreno} & string[][] & Matriz \code{filas $\times$ columnas} con nombres de terreno \\
\code{bases.A} & \code{\{fila, columna\}} & Posición base A \\
\code{bases.B} & \code{\{fila, columna\}} & Posición base B \\
\code{recurso} & \code{\{fila, columna\}} & Posición recurso (si nadie lo transporta) \\
\code{unidades} & array & Lista: \code{id, bando, tipo, fila, columna} \\
\code{turno} & \code{"A"} | \code{"B"} & Bando que debe actuar \\
\code{juego.portador\_recurso} & string | \code{null} & id unidad que porta, o \code{null} \\
\code{juego.profundidad\_maxima\_minimax} & int (opcional) & Profundidad para Minimax \\
\code{juego.limite\_turnos} & int (opcional) & Límite turnos (default 100) \\
\code{juego.variante} & string (opcional) & \code{base} o \code{dinamica} \\
\code{juego.turnos\_recuperacion} & int (opcional) & Recuperación propia de unidades dinámicas \\
\code{juego.inmunidad\_robo} & int (opcional) & Inmunidad del portador tras un robo \\
\code{utilidad} & object (opcional) & Pesos: entrega, recurso\_cargado, distancia\_objetivo, distancia\_rival \\
\code{prueba} & object (opcional) & Configuración de prueba concreta \\
\bottomrule
\caption{Campos del formato JSON (obligatorios + opcionales del equipo)}
\label{tab:campos-json}
\end{longtable}

\section{Normalización y Tolerancia}
\label{sec:normalizacion}

El loader (\code{scenario/loader.py}) acepta y normaliza:

\begin{itemize}
\item \textbf{Bandos:} \code{"A"}/\code{"a"} $\to$ \code{Side.A}, \code{"B"}/\code{"b"} $\to$ \code{Side.B}
\item \textbf{Costos:} enteros o decimales positivos (\code{float})
\item \textbf{Terrenos no transitables:} \code{costo} opcional (se ignora si \code{transitable: false})
\item \textbf{Campos propios:} \code{juego.limite\_turnos}, \code{utilidad}, \code{prueba.*} (RNF-02)
\item \textbf{Variante dinámica:} \code{juego.variante}, \code{juego.turnos\_recuperacion},
  \code{juego.inmunidad\_robo}; si no se indica, se usa \code{base}
\end{itemize}

\section{Ejemplo Oficial (Apéndice A)}
\label{sec:ejemplo-apendice-a}

\begin{minted}{json}
{
  "version": "1.0",
  "mapa": { "filas": 4, "columnas": 4 },
  "tipos_terreno": {
    "camino":  {"costo": 1, "transitable": true},
    "pasto":   {"costo": 2, "transitable": true},
    "bosque":  {"costo": 4, "transitable": true},
    "pantano": {"costo": 7, "transitable": true},
    "muro":    {"transitable": false}
  },
  "terreno": [
    ["camino", "camino", "pasto",  "pasto"],
    ["camino", "muro",   "bosque", "pasto"],
    ["camino", "camino", "pantano","pasto"],
    ["pasto",  "pasto",  "camino", "camino"]
  ],
  "bases": { "A": {"fila": 0, "columna": 0}, "B": {"fila": 3, "columna": 3} },
  "recurso": {"fila": 2, "columna": 2},
  "unidades": [
    {"id": "A1", "bando": "A", "tipo": "estandar", "fila": 0, "columna": 1},
    {"id": "A2", "bando": "A", "tipo": "estandar", "fila": 1, "columna": 0},
    {"id": "B1", "bando": "B", "tipo": "estandar", "fila": 3, "columna": 2},
    {"id": "B2", "bando": "B", "tipo": "estandar", "fila": 2, "columna": 3}
  ],
  "turno": "A",
  "juego": { "portador_recurso": null, "profundidad_maxima_minimax": 4 },
  "prueba": { "modo": "busqueda", "unidad_inicio": "A1", "objetivo": {"fila": 3, "columna": 2} }
}
\end{minted}

\section{Validación Atómica (RF-34, RF-35)}
\label{sec:validacion-atomica}

\code{load\_scenario(path)} implementa una validación \textbf{atómica}:

\begin{enumerate}
\item Lee y parsea JSON completo
\item Ejecuta \textbf{todas} las validaciones, acumulando errores
\item Si hay \textbf{cualquier} error $\to$ lanza \code{ScenarioError} con \textbf{lista completa}
  de mensajes; \textbf{nada} se construye ni modifica
\item Si todo OK $\to$ construye \code{Scenario(board, initial\_state, test\_config)}
\end{enumerate}

\section{9 Validaciones Mínimas (RF-35) + 1 Extra}
\label{sec:9-validaciones}

\begin{longtable}{p{1cm}p{3cm}p{10cm}}
\toprule
\textbf{\#} & \textbf{Validación} & \textbf{Código / Detalle} \\
\midrule
\endfirsthead
\toprule
\textbf{\#} & \textbf{Validación} & \textbf{Código / Detalle} \\
\midrule
\endhead
\midrule
\multicolumn{3}{r}{\textit{Continúa}} \\
\endfoot
\bottomrule
\endlastfoot
1 & Filas/columnas positivas y consistentes con matriz & \code{\_read\_dimensions} + \code{\_read\_grid} \\
2 & Terrenos usados existen en \code{tipos\_terreno} & \code{\_read\_terrains} + \code{\_read\_grid} \\
3 & Transitables tienen costo positivo & \code{\_read\_terrains} \\
4 & Bases, recurso, unidades dentro del mapa & \code{\_read\_bases}, \code{\_read\_position}, \code{\_read\_units} \\
5 & Ninguna unidad en celda no transitable & \code{\_read\_units} \\
6 & IDs de unidad únicos & \code{\_read\_units} \\
7 & Cada unidad en bando A/a o B/b & \code{\_read\_side} \\
8 & Turno es A/a o B/b & \code{\_read\_side} \\
9 & \code{portador\_recurso} es id existente o null & \code{\_read\_carrier} \\
+ & \textbf{Extra:} Ninguna unidad comparte celda & \code{\_read\_units} (mapa \code{cells}) \\
\bottomrule
\caption{10 validaciones (9 RF-35 + 1 extra del equipo)}
\label{tab:10-validaciones}
\end{longtable}

\section{Campos de \code{prueba} (G10, Decisión 2026-10-17)}
\label{sec:campos-prueba}

El enunciado no fija nombres canónicos para algoritmo, heurística, $k$, profundidad, orden.
El equipo definió campos opcionales propios en \code{prueba} (RNF-02), CLI sobrepone JSON:

\begin{table}[H]
\centering
\caption{Campos opcionales en \code{prueba}}
\begin{tabular}{lll}
\toprule
\textbf{Campo} & \textbf{Valores} & \textbf{Descripción} \\
\midrule
\code{modo} & \code{busqueda}, \code{adversarial}, \code{partida} & Modo de ejecución \\
\code{algoritmo} & \code{BFS}, \code{DFS}, \code{UCS}, \code{A\_STAR}, \code{BEAM}, \code{MINIMAX}, \code{ALPHA\_BETA} & Algoritmo a ejecutar \\
\code{heuristica} & \code{H1}, \code{H2} (\code{MANHATTAN}, \code{OBSTACLE\_AWARE}) & Heurística para A*/Beam \\
\code{k} & int $\ge$1 & Ancho de haz (Beam) \\
\code{profundidad} / \code{depth} & int $\ge$1 & Profundidad Minimax/$\alpha$-$\beta$ \\
\code{orden\_acciones} & \code{default}, \code{reversed} & Orden de sucesores \\
\code{unidad\_inicio} & string & id unidad de inicio (búsqueda) \\
\code{objetivo} & \code{\{fila, columna\}} & Meta de búsqueda \\
\bottomrule
\end{tabular}
\end{table}
\label{tab:prueba-campos}


Aliases canónicos: \code{A*}/\code{A\_ESTRELLA} $\to$ \code{A\_STAR}, \code{H1}/\code{MANHATTAN},
\code{H2}/\code{OBSTACLE\_AWARE}. CLI (\code{--algorithm}, \code{--heuristic}, \code{--k},
\code{--depth}, \code{--order}) \textbf{sobrepone} al JSON (G10).

\section{Ejemplos de Escenarios del Equipo}
\label{sec:escenarios-equipo}

\begin{itemize}
\item \code{apendice\_a.json}: Ejemplo oficial 4$\times$4
\item \code{caso\_01\_bfs\_vs\_dfs.json} -- Caso 1
\item \code{caso\_02\_bfs\_vs\_ucs.json} -- Caso 2
\item \code{caso\_03\_obstaculo\_heuristica.json} -- Caso 3
\item \code{caso\_04\_comparacion\_heuristicas.json} -- Caso 4
\item \code{caso\_05\_beam\_k.json} -- Caso 5
\item \code{caso\_07\_decision\_adversarial.json} -- Casos 6--10 (mapa 5$\times$7)
\item \code{prueba\_15x15.json} -- Rendimiento 15$\times$15 (3 unidades/bando)
\item \code{prueba\_20x20.json} -- Prueba 20$\times$20
\item \code{escenario\_03x03} .. \code{escenario\_30x30} -- Escalabilidad
\end{itemize}